% JAWS 2 @ ASE 2026 — Idea Paper (≤8 pages excl. references)
% ACM Primary Article Template, sigconf, single-anonymous
\PassOptionsToPackage{expansion=false}{microtype}
\documentclass[sigconf,nonacm]{acmart}

% JAWS is single-anonymous: author names visible
\settopmatter{authorsperrow=3}

\usepackage{booktabs}
\usepackage{multirow}
\usepackage{xcolor}

% Blind-review off (single-anonymous)
\setcopyright{none}
\renewcommand\footnotetextcopyrightpermission[1]{}
\pagestyle{plain}

\title{From Symptom to Structure: Analyzing Disagreement\\in Agent-Skill Security Scanners}






\begin{abstract}
Automated security analyses routinely disagree, yet evaluations treat that disagreement as a statistic---Cohen's $\kappa$, pairwise overlap---and stop there, as if it were noise or merely an outcome about tool quality.
We propose to treat disagreement differently: as \emph{analytical evidence about the analyses themselves}.
The structure of disagreement, once made comparable, exposes each analysis's coverage of the threat space and how it operationalizes detection there.
We develop a methodology for this conversion: a shared observation language that makes disagreement localizable, a decomposition into \emph{coverage gaps} and \emph{operational divergence}, blind-spot hypotheses locked before construction, and adversarial validation.
We instantiate the idea on agent-skill security scanners---an emerging supply-chain defense whose scanners disagree sharply.
Across 582 real malicious skills and three heterogeneous scanners, the decomposition predicts concrete blind spots, and locked adversarial construction confirms them: hidden-file payloads evade SkillSpector on 14/15, expression variants evade Cisco on 6/6, and combined strategies evade SkillSpector on 5/5, with selected failures traced to source-level mechanisms.
Disagreement, we argue, is not merely a statistic to report---it is evidence to mine.
\end{abstract}

\keywords{agent skills, security scanners, taxonomy, blind spots, adversarial construction}

\begin{document}
\maketitle

% ============================================================
\section{Introduction}
\label{sec:intro}

Automated security analyses disagree.
Given the same artifact, one scanner flags it, another passes it, and a third flags something else entirely.
The standard response is to \emph{measure} the disagreement---Cohen's $\kappa$, pairwise overlap, per-tool recall---and treat the result as an evaluation outcome: a fact about tool quality, or noise to be averaged away.
We ask a different question: \emph{what does the pattern of disagreement reveal about the analyses themselves?}
We argue that disagreement is not merely a statistic to report.
It is analytical evidence: once disagreements are made comparable, their structure exposes what each analysis covers, what it cannot see, and why.

Today this evidence is inaccessible, because analyses speak incompatible languages.
One scanner reports ``tool poisoning,'' another ``supply-chain risk,'' a third a family of regex matches.
There is no shared frame in which to ask \emph{where exactly} they disagree, so disagreement remains global---a $\kappa$---but never localizable.
Without localization there is no explanation; without explanation, disagreement cannot be acted upon.

We propose to convert disagreement into structure.
The conversion has four steps.
First, a \emph{shared observation language}---a coordinate space over the threat space---into which each scanner's categories are mapped, making disagreement localizable.
Second, a \emph{structural decomposition}: localized disagreement separates into \emph{coverage gaps} (a scanner does not inspect a region at all) and \emph{operational divergence} (scanners cover the same region but instantiate detection differently).
Third, the decomposition yields \emph{blind-spot hypotheses}, which we lock before building anything.
Fourth, \emph{adversarial validation}: coordinate-driven construction tests the locked hypotheses, and selected failures are traced to implementation mechanisms.

We instantiate this idea on \emph{agent-skill security scanners}.
Coding agents extend their capabilities by installing skills---packages of instructions and scripts---and malicious skills are a live supply-chain threat (the ClawHavoc incident surfaced over 1{,}184 of them).
Scanners have emerged to screen skills, but they disagree sharply, and no existing evaluation explains the disagreement.
The object is timely; the methodology is the contribution.

An honesty point frames our results.
In aggregate the scanners are strong: across 582 real malicious skills, SkillSpector detects 92.8\%, Cisco 85.2\%, Caterpillar 72.5\%.
The blind spots are therefore invisible in average recall; they emerge only when the threat space is partitioned.
This is precisely the structure that aggregate metrics conceal---and precisely what disagreement, treated as evidence, reveals.

\noindent\textbf{Contributions.} This is an idea paper: one general claim, a methodology, an instantiation, and early evidence.
\begin{itemize}
  \item \textbf{Idea.} Disagreement among automated security analyses is analytical evidence, not merely an evaluation outcome.
  \item \textbf{Methodology.} A conversion pipeline---shared observation language $\rightarrow$ structural decomposition $\rightarrow$ locked hypotheses $\rightarrow$ adversarial validation (\S\ref{sec:measure}--\S\ref{sec:validate}).
  \item \textbf{Instantiation.} Agent-skill security scanners, studied through a three-dimensional attack-coordinate space.
  \item \textbf{Early evidence.} 582 real malicious skills, three scanners, 129 locked adversarial constructions, and source-level mechanism tracing.
\end{itemize}

\noindent\textbf{Scope.} We demonstrate the loop end-to-end on one instantiation.
The generality of the idea---to package scanners, dependency scanners, static analyzers, LLM-based security evaluators---is the research direction this early result opens, to be developed in an extended version.

\noindent\textbf{Results.} Hidden-file payloads evade SkillSpector on 14/15 constructed samples (its build context skips dotfiles); expression-layer variants evade Cisco on 6/6 (its pipeline requires literal command surfaces); combined strategies evade SkillSpector on 5/5.
A regex-only reference tool is blind to all 10 semantic variants, showing blind spots deepen as detection grows more primitive.

\noindent\textbf{Paper identity.} Disagreement is not merely a number: it reveals structure that can be hypothesized about and tested.

% ============================================================
\section{Background and Motivation}
\label{sec:background}

Agent skills are distributed through marketplaces such as ClawHub and skills.sh. Because a skill can carry arbitrary instructions and scripts, it is an active supply-chain attack surface: the ClawHavoc incident identified over 1{,}184 malicious skills in a single marketplace. Security scanners have emerged to screen skills before installation.

We study three scanners spanning the mainstream detection architectures. \textbf{SkillSpector} (NVIDIA) combines static rules, AST analysis, YARA signatures, and an LLM semantic layer. \textbf{Cisco Skill Scanner} (AI Defense) combines static analysis, YARA, LLM-based judgment, and behavioral dataflow. \textbf{Caterpillar} is a lightweight regex-based tool. SkillSpector and Cisco are production-grade; Caterpillar serves as a lower-bound reference showing how blind spots deepen as detection grows more primitive.

These scanners frequently disagree on the same skill. Recent measurements quantify the disagreement: across 67{,}453 skills, pairwise overlap of flagged sets is at most 10.4\%, and only 0.69\% are flagged by all scanners. Yet such numbers stay descriptive: they tell us \emph{that} scanners disagree, not \emph{where} in the attack space the disagreement arises or \emph{why}.

% ============================================================
\section{A Shared Observation Language}
\label{sec:measure}

The reason disagreement cannot currently be interrogated is that the scanners speak incompatible languages: SkillSpector speaks of 17 risk classes, Cisco of 18 threat categories, Caterpillar of regex rule families. There is no frame in which to ask \emph{where exactly} they disagree. The coordinate space we introduce is, first and foremost, a \emph{shared observation language}---the instrument that makes disagreement localizable. It is not a universal taxonomy; it is the observation layer required to interrogate disagreement.

\subsection{Three Dimensions}
Each threat is located by three dimensions: \emph{source} (where it enters: supply chain, user input, external untrusted content, runtime environment, or unknown), \emph{mechanism} (how it attacks: code execution, instruction manipulation, dependency manipulation, privilege abuse, obfuscation, state corruption, \ldots), and \emph{target} (what it seeks: information theft, persistent control, resource abuse, system damage, financial theft, content-safety harm, defense evasion). The axes follow established security modeling: the attacker model (cf.\ Dolev--Yao), the technique axis (cf.\ ATT\&CK), and the security-property axis (cf.\ the CIA triad).

\subsection{Mapping Scanner Languages}
We mapped 67 scanner-native categories into the space, yielding 43 coordinates. To check that the space can carry existing threat descriptions, we mapped 1{,}253 threat statements from 84 sources. Only 17.9\% map cleanly to a full triple; \textbf{56.5\% are partial triples}---existing vocabularies cannot even fully express them---and 20.3\% fall into genuine gaps. This partial-triple rate is our key measurement-space evidence: current categorizations do not provide a complete description; a coordinate space does.

\subsection{Orthogonality}
The dimensions are near-orthogonal (normalized mutual information between source and target is 0.232), and ablating any single axis collapses coordinate resolution. The space therefore carries non-redundant information and serves as the observation layer for all subsequent analysis.

% ============================================================
\section{Structural Disagreement Analysis}
\label{sec:explain}

With a common observation layer in place, we project each scanner's detections onto the coordinate space and compare scanners coordinate by coordinate. This is the paper's core analysis, measured on 582 real malicious skills (350 wild samples stratified by behavior, 232 coordinate-generated).

\subsection{Aggregate Detection Conceals the Structure}
In aggregate the scanners are strong: SkillSpector detects 92.8\%, Cisco 85.2\%, Caterpillar 72.5\%. Stopping here, one would conclude detection is largely solved. The blind spots are invisible at this level of abstraction; they emerge only when the attack space is partitioned.

\subsection{Coverage Gaps}
The first structural cause of disagreement is a \emph{coverage gap}: a scanner does not inspect a coordinate at all. The source axis is collectively under-determined---after mapping to coordinates, inter-scanner agreement on source is negative ($\kappa_{source}=-0.137$), because scanners report behavior, not provenance. Coordinates whose source is runtime environment or user input are systematically less covered than supply-chain coordinates.

\subsection{Operational Divergence: Semantic Alignment, Operational Divergence}
The second cause is subtler. Scanners can \emph{semantically agree} that a coordinate matters while \emph{operationally diverging} on how to detect it. Consider the coordinate (credential access $\times$ code execution $\times$ data exfiltration): all three scanners nominally cover it, yet SkillSpector instantiates detection through script-level semantic evidence, Cisco through literal command triggers, and Caterpillar through lexical regex patterns. Same coordinate, same semantic intent, three different operational surfaces. An attack that removes the literal surface evades Cisco and Caterpillar while SkillSpector may still catch it---and vice versa for attacks that hide the script.

\subsection{Misclassification}
A third, orthogonal finding: scanners can detect a behavior but misclassify its intent. In a pilot audit of 30 SkillSpector findings, every one described a real capability that was in fact a legitimate, declared tool behavior---``dangerous but legal.'' Detection fired; the verdict was wrong.

% ============================================================
\section{From Analysis to Blind-Spot Hypotheses}
\label{sec:hypotheses}

The coordinate space is not itself a blind-spot oracle. It is a common observation layer through which two kinds of evidence are turned into explicit, pre-specified hypotheses, which we lock \emph{before} building any adversarial sample. To prevent post-hoc rationalization, we separate exploratory experiments from confirmatory ones.

\subsection{Two Evidence Sources}
\emph{Type A (coordinate-coverage evidence).} A coordinate uncovered or weakly covered in the coverage matrix yields the hypothesis that the scanner will miss attacks occupying that coordinate.

\emph{Type B (scanner-architecture evidence).} An implementation property yields a mechanistic limitation. Reading SkillSpector's source, for instance, reveals that its build context skips dotfiles; this yields the hypothesis that hidden-file payloads will evade analysis.

\subsection{Hypothesis-Locking Protocol}
For each hypothesis we record the evidence, the locked prediction, the construction specification, and the validation outcome. Locking precedes construction: the construction specification is written from the hypothesis, not from observed misses. Our timeline is explicit: the coverage matrix and blind-spot hypotheses were formed \emph{before} the confirmatory construction experiments (Arm7 onward); earlier arms (Arm1--6) are labeled exploratory and serve only to establish that coordinate-driven construction is feasible.

% ============================================================
\section{Early Evidence: Adversarial Validation}
\label{sec:validate}

We now demonstrate that the loop works end-to-end. This is early evidence---proof of feasibility for the idea and the methodology---rather than a claim of benchmark completeness. We validate locked hypotheses through coordinate-driven adversarial construction: 129 constructed samples across 13 experiment arms, each scanned by all three scanners.

\subsection{Centerpiece: Hidden-File Payloads}
\emph{Evidence (Type B):} SkillSpector's build context walks only non-dotfiles.
\emph{Locked hypothesis:} payloads placed in hidden files (.env, .profile, .gitconfig) will evade SkillSpector.
\emph{Construction:} 15 skills whose SKILL.md is benign while the payload lives entirely in a dotfile.
\emph{Outcome:} \textbf{14/15 evade} SkillSpector (score $\le$50); 11 score 0. The mechanism is confirmed at source level: the dotfiles never enter the scanned context.

\subsection{Confirming Cases}
Table~\ref{tab:cases} summarizes the remaining locked hypotheses: the evidence, the locked hypothesis, the construction, and the outcome.

\begin{table}[t]
\caption{Locked blind-spot hypotheses and validation outcomes.}
\label{tab:cases}
\small
\begin{tabular}{@{}p{1.1cm}p{2.6cm}p{3.4cm}p{2.9cm}@{}}
\toprule
Scanner & Evidence & Locked hypothesis & Outcome \\
\midrule
Cisco & literal-trigger dependence & expression variants lack a detection surface & 0/6 detected \\
Cat & regex-only matching & semantic variants without lexical markers evade & 0/10 detected \\
SS & multi-strategy interaction & combined mechanisms escape isolated checks & 5/5 bypass \\
SS & wild variable-injection & pipeline injection evades semantic analysis & 3/5 score 0 \\
\bottomrule
\end{tabular}
\end{table}

\subsection{Honest Scoping}
Two caveats bound our claims. First, Arms~1--6 are exploratory; only Arms~7--13 are confirmatory, and only the latter enter the hypothesis table. Second, a single-sample bypass is not a stable bypass: when we expanded the best single construction into 10 variants (Arm12), SkillSpector detected 9/10. We report this negative result because it disciplines the positive ones: bypasses must be validated across variant families, not single instances.

% ============================================================
\section{Implications and Future Work}
\label{sec:implications}

\emph{For scanner evaluation.} Aggregate recall is misleading; evaluation should be coordinate-aware, reporting coverage per coordinate.

\emph{For the skill ecosystem.} Single-scanner governance inherits that scanner's blind spots wholesale. Defense-in-depth across scanners with complementary operational surfaces is warranted.

\emph{For adversarial methodology.} Hypothesis locking separates genuine prediction from post-hoc rationalization and should become standard practice.

\emph{Future work.} We are extending this into a full benchmark and a taxonomy-guided defensive filter, to be reported in an extended version.

% ============================================================
\section{Related Work}
\label{sec:related}

Prior work addresses individual parts of this problem. Malicious-skill benchmarks (MalSkillBench, SkillTrustBench) provide corpora but not a shared cross-scanner measurement space. Scanner evaluations (ClawHub, Locate-and-Judge, SkillsMetric) quantify disagreement or recall but do not structurally attribute it. Adversarial-generation studies expose failures but are not driven by an explicit attack-space hypothesis. Implementation analyses identify individual bugs but do not connect them to a cross-scanner attack model. Our contribution connects these steps into a single analysis loop.

Adversarial evasion of security tools is well studied in malware detection (e.g., semantics-preserving transformations). Our setting differs in that the artifact is a hybrid of prose instructions and code, and the blind spots we find are architectural rather than feature-space perturbations.

% ============================================================
\bibliographystyle{ACM-Reference-Format}
\bibliography{references}

\end{document}
